\chapter{Conclusion, Limitations and Future Work}\label{ch:concl}

\section{Conclusion}
This project built Food Price Watch, a monitoring and abnormal-rise alert system for the TCB daily retail prices of eight food items in Dhaka, and evaluated it with every setting chosen on 2024 and committed before 2025-26 was scored. The prices are sticky: an item's price is unchanged on 62 to 87\,\% of report days, yet sustained 10\,\% rises do happen (70 in 2021-26), mostly in onion, potato and imported garlic (RQ1). Neither ETS nor LightGBM beats \naive{} at any horizon, in 2024 (relative MAE 1.00 to 1.11) or in 2025-26 \retro{} (1.01 to 1.21), and every significant difference favours \naive{} (RQ2). The alerts do catch sustained rises: the ML range caught 6 of 11 in 2024 and 17 of 23 in 2025-26 \retro{}, and the simple range 8 of 11 and 15 of 23, about 2.5 to 2.8 times the recall of the same alerts on random dates ($p = 0.003$), a median of 1.5 to 2 report days after a rise begins (RQ3). The two ranges catch nearly the same rises, so a daily monitor for these prices does not need machine learning: a range drawn from each item's own past price changes is enough (RQ4).

\section{Limitations}
\begin{itemize}[leftmargin=*,itemsep=2pt]
\item The 2025-26 results are retrospective: that period had been looked at in an earlier run of the project.
\item The eight items were chosen after the earlier results had been seen.
\item There are few sustained rises (11 in 2024, 23 in 2025-26), and none for three of the items.
\item Rule C is close to the definition of a sustained rise, so its score is partly circular.
\item The range was slightly too narrow in 2025-26 (coverage 93.4 to 93.9\,\% against 95\,\%).
\item Only rises are monitored; only Dhaka is covered; and the monitor replays past data rather than receiving new sheets.
\item The price is the midpoint of TCB's lowest and highest price, and its accuracy depends on TCB's own recording.
\end{itemize}

\section{Future Work}
\begin{itemize}[leftmargin=*,itemsep=2pt]
\item Connect the monitor to new TCB sheets as they are published, with TCB's permission, and run it forward on genuinely unseen days.
\item Extend it to more items, to falls as well as rises, and to other cities.
\item Make the range adapt faster when an item becomes more volatile, to bring coverage back to its nominal level.
\item Compare the alerts with FAO's and WFP's indicators on the same items, and test whether weekly or monthly drivers help at longer horizons.
\end{itemize}
